\chapter{Complexity — IS Count Regs}
%%% generated by the blueprint pipeline from TCSlib.Complexity.SpaceComplexity.Logspace.ISCountRegs
\section{Overview}

Declarations of the Lean module \texttt{TCSlib.Complexity.SpaceComplexity.Logspace.ISCountRegs}.

\section{Declarations}

\begin{definition}[First register index]
\lean{Complexity.LogProg.ISCount.rn}
\label{Complexity.LogProg.ISCount.rn}
\leanok
The first register index among seventeen registers is $0$.
\end{definition}

\begin{definition}[Second register index]
\lean{Complexity.LogProg.ISCount.rlg}
\label{Complexity.LogProg.ISCount.rlg}
\leanok
The second register index among seventeen registers is $1$.
\end{definition}

\begin{definition}[Rank coordinate]
\lean{Complexity.LogProg.ISCount.rk}
\label{Complexity.LogProg.ISCount.rk}
\leanok
The coordinate is the element $2$ of the finite set $\{0,1,\ldots,16\}$, regarded as an
element of $\mathrm{Fin}(17)$.
\end{definition}

\begin{definition}[Size coordinate]
\lean{Complexity.LogProg.ISCount.rs}
\label{Complexity.LogProg.ISCount.rs}
\leanok
The coordinate is the element $3$ of the finite set $\{0,1,\ldots,16\}$, regarded as an
element of $\mathrm{Fin}(17)$.
\end{definition}

\begin{definition}[Real coordinate]
\lean{Complexity.LogProg.ISCount.re}
\label{Complexity.LogProg.ISCount.re}
\leanok
The coordinate is the element $4$ of the finite set $\{0,1,\ldots,16\}$, regarded as an
element of $\mathrm{Fin}(17)$.
\end{definition}

\begin{definition}[Capital rank coordinate]
\lean{Complexity.LogProg.ISCount.rK}
\label{Complexity.LogProg.ISCount.rK}
\leanok
The coordinate is the element $5$ of the finite set $\{0,1,\ldots,16\}$, regarded as an
element of $\mathrm{Fin}(17)$.
\end{definition}

\begin{definition}[Total rank coordinate]
\lean{Complexity.LogProg.ISCount.rT}
\label{Complexity.LogProg.ISCount.rT}
\leanok
The coordinate is the element $6$ of the finite set $\{0,1,\ldots,16\}$, regarded as an
element of $\mathrm{Fin}(17)$.
\end{definition}

\begin{definition}[Index coordinate]
\lean{Complexity.LogProg.ISCount.ri}
\label{Complexity.LogProg.ISCount.ri}
\leanok
The coordinate is the element $7$ of the finite set $\{0,1,\ldots,16\}$, regarded as an
element of $\mathrm{Fin}(17)$.
\end{definition}

\begin{definition}[First index coordinate]
\lean{Complexity.LogProg.ISCount.ri1}
\label{Complexity.LogProg.ISCount.ri1}
\leanok
The coordinate is the element $8$ of the finite set $\{0,1,\ldots,16\}$, regarded as an
element of $\mathrm{Fin}(17)$.
\end{definition}

\begin{definition}[Complement coordinate]
\lean{Complexity.LogProg.ISCount.rcp}
\label{Complexity.LogProg.ISCount.rcp}
\leanok
The coordinate is the element $9$ of the finite set $\{0,1,\ldots,16\}$, regarded as an
element of $\mathrm{Fin}(17)$.
\end{definition}

\begin{definition}[Vertex register index]
\lean{Complexity.LogProg.ISCount.rv}
\label{Complexity.LogProg.ISCount.rv}
\leanok
Let the 17-element index set be $\{0,1,\ldots,16\}$. The designated index $rv$ is $10$.
\end{definition}

\begin{definition}[Code register index]
\lean{Complexity.LogProg.ISCount.rc}
\label{Complexity.LogProg.ISCount.rc}
\leanok
Let the 17-element index set be $\{0,1,\ldots,16\}$. The designated index $rc$ is $11$.
\end{definition}

\begin{definition}[Update register index]
\lean{Complexity.LogProg.ISCount.ru}
\label{Complexity.LogProg.ISCount.ru}
\leanok
Let the 17-element index set be $\{0,1,\ldots,16\}$. The designated index $ru$ is $12$.
\end{definition}

\begin{definition}[Count register index]
\lean{Complexity.LogProg.ISCount.rcnt}
\label{Complexity.LogProg.ISCount.rcnt}
\leanok
Let the 17-element index set be $\{0,1,\ldots,16\}$. The designated index $rcnt$ is
$13$.
\end{definition}

\begin{definition}[Guess register index]
\lean{Complexity.LogProg.ISCount.rg}
\label{Complexity.LogProg.ISCount.rg}
\leanok
Let the 17-element index set be $\{0,1,\ldots,16\}$. The designated index $rg$ is $14$.
\end{definition}

\begin{definition}[Current register index]
\lean{Complexity.LogProg.ISCount.rcur}
\label{Complexity.LogProg.ISCount.rcur}
\leanok
Let the 17-element index set be $\{0,1,\ldots,16\}$. The designated index $rcur$ is
$15$.
\end{definition}

\begin{definition}[Step register index]
\lean{Complexity.LogProg.ISCount.rstp}
\label{Complexity.LogProg.ISCount.rstp}
\leanok
Let the 17-element index set be $\{0,1,\ldots,16\}$. The designated index $rstp$ is
$16$.
\end{definition}

\begin{definition}[Counting algorithm state]
\lean{Complexity.LogProg.ISCount.St}
\label{Complexity.LogProg.ISCount.St}
\leanok
A state is a tuple of natural numbers recording the input length $n$, halving count
$lg$, add-twice counter $k$, window radius $s$, exponent $e$ and code bound $K$,
doubling scratch register $T$, stage $i$ and stage-plus-one $i1$, known-ball size $cp$,
outer vertex $v$ and next-ball count $c$, inner vertex $u$ and inner count $cnt$,
guessed code $g$, current walk code $cur$, and walk step count $stp$.
\end{definition}

\begin{definition}[Register vector]
\lean{Complexity.LogProg.ISCount.St.toV}
\label{Complexity.LogProg.ISCount.St.toV}
\leanok
\uses{Complexity.LogProg.ISCount.St}
A register state consists of 17 natural-number registers, and its register vector lists,
in order, the input length, halving count, add-twice counter, window radius, code-bound
exponent, code bound, doubling scratch register, stage, stage plus one, known-ball size,
outer vertex, next-ball count, inner vertex, inner count, guessed code, current walk
code, and walk step count.
\end{definition}

\begin{definition}[Zero register state]
\lean{Complexity.LogProg.ISCount.St.zero}
\label{Complexity.LogProg.ISCount.St.zero}
\leanok
\uses{Complexity.LogProg.ISCount.St}
The zero register state is the state in which all 17 registers are zero.
\end{definition}

\begin{lemma}[Zero register vector]
\lean{Complexity.LogProg.ISCount.St.toV_zero}
\label{Complexity.LogProg.ISCount.St.toV_zero}
\leanok
\uses{Complexity.LogProg.ISCount.St.toV, Complexity.LogProg.ISCount.St.zero}
\difficulty{4}
The register vector of the state whose 17 natural-number fields are all zero is the
constant function on the 17 register indices with value $0$.
\end{lemma}

\begin{lemma}[Input-length coordinate]
\lean{Complexity.LogProg.ISCount.St.toV_n}
\label{Complexity.LogProg.ISCount.St.toV_n}
\leanok
\uses{Complexity.LogProg.ISCount.St.toV, Complexity.LogProg.ISCount.rn}
\difficulty{1}
For every counting-algorithm state $\sigma$ and natural number $z$, the register vector
of $\sigma$ has value $n$ at coordinate $0$.
\end{lemma}

\begin{lemma}[Halving-count coordinate]
\lean{Complexity.LogProg.ISCount.St.toV_lg}
\label{Complexity.LogProg.ISCount.St.toV_lg}
\leanok
\uses{Complexity.LogProg.ISCount.St.toV, Complexity.LogProg.ISCount.rlg}
\difficulty{1}
For every counting-algorithm state $\sigma$ and natural number $z$, the register vector
of $\sigma$ has value $lg$ at coordinate $1$.
\end{lemma}

\begin{lemma}[Add-twice counter coordinate]
\lean{Complexity.LogProg.ISCount.St.toV_k}
\label{Complexity.LogProg.ISCount.St.toV_k}
\leanok
\uses{Complexity.LogProg.ISCount.St.toV, Complexity.LogProg.ISCount.rk}
\difficulty{1}
For every counting-algorithm state $\sigma$ and natural number $z$, the register vector
of $\sigma$ has value $k$ at coordinate $2$.
\end{lemma}

\begin{lemma}[Window-radius coordinate]
\lean{Complexity.LogProg.ISCount.St.toV_s}
\label{Complexity.LogProg.ISCount.St.toV_s}
\leanok
\uses{Complexity.LogProg.ISCount.St.toV, Complexity.LogProg.ISCount.rs}
\difficulty{1}
For every counting-algorithm state $\sigma$ and natural number $z$, the register vector
of $\sigma$ has value $s$ at coordinate $3$.
\end{lemma}

\begin{lemma}[Code-exponent coordinate]
\lean{Complexity.LogProg.ISCount.St.toV_e}
\label{Complexity.LogProg.ISCount.St.toV_e}
\leanok
\uses{Complexity.LogProg.ISCount.St.toV, Complexity.LogProg.ISCount.re}
\difficulty{1}
For every counting-algorithm state $\sigma$ and natural number $z$, the register vector
of $\sigma$ has value $\sigma.e$ at coordinate $4$.
\end{lemma}

\begin{lemma}[Code-bound register coordinate]
\lean{Complexity.LogProg.ISCount.St.toV_K}
\label{Complexity.LogProg.ISCount.St.toV_K}
\leanok
\uses{Complexity.LogProg.ISCount.St.toV, Complexity.LogProg.ISCount.rK}
\difficulty{1}
For every counting-algorithm state $\sigma$ and natural number $z$, the register vector
of $\sigma$ has value $\sigma.K$ at coordinate $5$.
\end{lemma}

\begin{lemma}[Scratch-register coordinate]
\lean{Complexity.LogProg.ISCount.St.toV_T}
\label{Complexity.LogProg.ISCount.St.toV_T}
\leanok
\uses{Complexity.LogProg.ISCount.St.toV, Complexity.LogProg.ISCount.rT}
\difficulty{1}
For every counting-algorithm state $\sigma$ and natural number $z$, the register vector
of $\sigma$ has value $\sigma.T$ at coordinate $6$.
\end{lemma}

\begin{lemma}[Stage coordinate]
\lean{Complexity.LogProg.ISCount.St.toV_i}
\label{Complexity.LogProg.ISCount.St.toV_i}
\leanok
\uses{Complexity.LogProg.ISCount.St.toV, Complexity.LogProg.ISCount.ri}
\difficulty{1}
For every counting-algorithm state $\sigma$ and natural number $z$, the register vector
of $\sigma$ has value $i$ at coordinate $7$.
\end{lemma}

\begin{lemma}[Stage-plus-one coordinate]
\lean{Complexity.LogProg.ISCount.St.toV_i1}
\label{Complexity.LogProg.ISCount.St.toV_i1}
\leanok
\uses{Complexity.LogProg.ISCount.St.toV, Complexity.LogProg.ISCount.ri1}
\difficulty{1}
For every counting-algorithm state $\sigma$ and natural number $z$, the register vector
of $\sigma$ has value $i1$ at coordinate $8$.
\end{lemma}

\begin{lemma}[Known-ball size coordinate]
\lean{Complexity.LogProg.ISCount.St.toV_cp}
\label{Complexity.LogProg.ISCount.St.toV_cp}
\leanok
\uses{Complexity.LogProg.ISCount.St.toV, Complexity.LogProg.ISCount.rcp}
\difficulty{1}
For every counting-algorithm state $\sigma$ and natural number $z$, the register vector
of $\sigma$ has value $\sigma.\mathrm{cp}$ at coordinate $9$.
\end{lemma}

\begin{lemma}[Outer-vertex coordinate]
\lean{Complexity.LogProg.ISCount.St.toV_v}
\label{Complexity.LogProg.ISCount.St.toV_v}
\leanok
\uses{Complexity.LogProg.ISCount.St.toV, Complexity.LogProg.ISCount.rv}
\difficulty{1}
For any counting-algorithm state $\sigma$, its register vector has value $\sigma.v$ at
index $10$.
\end{lemma}

\begin{lemma}[Next-ball count coordinate]
\lean{Complexity.LogProg.ISCount.St.toV_c}
\label{Complexity.LogProg.ISCount.St.toV_c}
\leanok
\uses{Complexity.LogProg.ISCount.St.toV, Complexity.LogProg.ISCount.rc}
\difficulty{1}
For every counting-algorithm state $\sigma$ and natural number $z$, the register vector
of $\sigma$ has value $\sigma.c$ at coordinate $11$.
\end{lemma}

\begin{lemma}[Inner-vertex coordinate]
\lean{Complexity.LogProg.ISCount.St.toV_u}
\label{Complexity.LogProg.ISCount.St.toV_u}
\leanok
\uses{Complexity.LogProg.ISCount.St.toV, Complexity.LogProg.ISCount.ru}
\difficulty{1}
For any counting-algorithm state $\sigma$, its register vector has value $\sigma.u$ at
index $12$.
\end{lemma}

\begin{lemma}[Inner-count coordinate]
\lean{Complexity.LogProg.ISCount.St.toV_cnt}
\label{Complexity.LogProg.ISCount.St.toV_cnt}
\leanok
\uses{Complexity.LogProg.ISCount.St.toV, Complexity.LogProg.ISCount.rcnt}
\difficulty{1}
For every counting-algorithm state $\sigma$ and natural number $z$, the register vector
of $\sigma$ has value $\sigma.\mathrm{cnt}$ at coordinate $13$.
\end{lemma}

\begin{lemma}[Guessed code coordinate]
\lean{Complexity.LogProg.ISCount.St.toV_g}
\label{Complexity.LogProg.ISCount.St.toV_g}
\leanok
\uses{Complexity.LogProg.ISCount.St.toV, Complexity.LogProg.ISCount.rg}
\difficulty{1}
For every counting-algorithm state $\sigma$ and natural number $z$, the register vector
of $\sigma$ has value $g$ at coordinate $14$.
\end{lemma}

\begin{lemma}[Current walk-code coordinate]
\lean{Complexity.LogProg.ISCount.St.toV_cur}
\label{Complexity.LogProg.ISCount.St.toV_cur}
\leanok
\uses{Complexity.LogProg.ISCount.St.toV, Complexity.LogProg.ISCount.rcur}
\difficulty{1}
For every counting-algorithm state $\sigma$ and natural number $z$, the register vector
of $\sigma$ has value $\sigma.\mathrm{cur}$ at coordinate $15$.
\end{lemma}

\begin{lemma}[Walk-step-count coordinate]
\lean{Complexity.LogProg.ISCount.St.toV_stp}
\label{Complexity.LogProg.ISCount.St.toV_stp}
\leanok
\uses{Complexity.LogProg.ISCount.St.toV, Complexity.LogProg.ISCount.rstp}
\difficulty{1}
For any counting-algorithm state $\sigma$, its register vector has value $\sigma.stp$ at
index $16$.
\end{lemma}

\begin{lemma}[Updating the input length]
\lean{Complexity.LogProg.ISCount.St.upd_n}
\label{Complexity.LogProg.ISCount.St.upd_n}
\leanok
\uses{Complexity.LogProg.ISCount.St.toV, Complexity.LogProg.ISCount.rn}
\difficulty{4}
For any counting-algorithm state $\sigma$ and natural number $z$, replacing coordinate
$0$ of the register vector of $\sigma$ by $z$ gives the register vector of the state
obtained by replacing its input-length register $n$ by $z$.
\end{lemma}

\begin{lemma}[Updating the halving count]
\lean{Complexity.LogProg.ISCount.St.upd_lg}
\label{Complexity.LogProg.ISCount.St.upd_lg}
\leanok
\uses{Complexity.LogProg.ISCount.St.toV, Complexity.LogProg.ISCount.rlg}
\difficulty{4}
For any counting-algorithm state $\sigma$ and natural number $z$, replacing coordinate
$1$ of the register vector of $\sigma$ by $z$ gives the register vector of the state
obtained by replacing its halving count $lg$ by $z$.
\end{lemma}

\begin{lemma}[Updating the add-twice counter]
\lean{Complexity.LogProg.ISCount.St.upd_k}
\label{Complexity.LogProg.ISCount.St.upd_k}
\leanok
\uses{Complexity.LogProg.ISCount.St.toV, Complexity.LogProg.ISCount.rk}
\difficulty{4}
For any counting-algorithm state $\sigma$ and natural number $z$, replacing coordinate
$2$ of the register vector of $\sigma$ by $z$ gives the register vector of the state
obtained by replacing its add-twice counter $k$ by $z$.
\end{lemma}

\begin{lemma}[Updating the window radius]
\lean{Complexity.LogProg.ISCount.St.upd_s}
\label{Complexity.LogProg.ISCount.St.upd_s}
\leanok
\uses{Complexity.LogProg.ISCount.St.toV, Complexity.LogProg.ISCount.rs}
\difficulty{4}
For any counting-algorithm state $\sigma$ and natural number $z$, replacing coordinate
$3$ of the register vector of $\sigma$ by $z$ gives the register vector of the state
obtained by replacing its window-radius register $s$ by $z$.
\end{lemma}

\begin{lemma}[Updating the code-bound exponent]
\lean{Complexity.LogProg.ISCount.St.upd_e}
\label{Complexity.LogProg.ISCount.St.upd_e}
\leanok
\uses{Complexity.LogProg.ISCount.St.toV, Complexity.LogProg.ISCount.re}
\difficulty{4}
For any counting-algorithm state $\sigma$ and natural number $z$, replacing coordinate
$4$ of its register vector by $z$ yields the register vector of the state obtained by
replacing its code-bound exponent $e$ by $z$.
\end{lemma}

\begin{lemma}[Updating the code-bound register]
\lean{Complexity.LogProg.ISCount.St.upd_K}
\label{Complexity.LogProg.ISCount.St.upd_K}
\leanok
\uses{Complexity.LogProg.ISCount.St.toV, Complexity.LogProg.ISCount.rK}
\difficulty{4}
For any counting-algorithm state $\sigma$ and natural number $z$, updating its register
vector at index $5$ to $z$ gives the register vector of the state obtained by replacing
its code bound $K$ with $z$.
\end{lemma}

\begin{lemma}[Updating the scratch register]
\lean{Complexity.LogProg.ISCount.St.upd_T}
\label{Complexity.LogProg.ISCount.St.upd_T}
\leanok
\uses{Complexity.LogProg.ISCount.St.toV, Complexity.LogProg.ISCount.rT}
\difficulty{4}
For any counting-algorithm state $\sigma$ and natural number $z$, updating its register
vector at index $6$ to $z$ gives the register vector of the state obtained by replacing
its doubling scratch register $T$ with $z$.
\end{lemma}

\begin{lemma}[Updating the stage coordinate]
\lean{Complexity.LogProg.ISCount.St.upd_i}
\label{Complexity.LogProg.ISCount.St.upd_i}
\leanok
\uses{Complexity.LogProg.ISCount.St.toV, Complexity.LogProg.ISCount.ri}
\difficulty{4}
For any counting-algorithm state $\sigma$ and natural number $z$, replacing coordinate
$7$ of its register vector by $z$ yields the register vector of the state obtained by
replacing its stage $i$ by $z$.
\end{lemma}

\begin{lemma}[Updating the next-stage coordinate]
\lean{Complexity.LogProg.ISCount.St.upd_i1}
\label{Complexity.LogProg.ISCount.St.upd_i1}
\leanok
\uses{Complexity.LogProg.ISCount.St.toV, Complexity.LogProg.ISCount.ri1}
\difficulty{4}
For any counting-algorithm state $\sigma$ and natural number $z$, replacing coordinate
$8$ of its register vector by $z$ yields the register vector of the state obtained by
replacing its stage-plus-one value $i1$ by $z$.
\end{lemma}

\begin{lemma}[Updating the known-ball size]
\lean{Complexity.LogProg.ISCount.St.upd_cp}
\label{Complexity.LogProg.ISCount.St.upd_cp}
\leanok
\uses{Complexity.LogProg.ISCount.St.toV, Complexity.LogProg.ISCount.rcp}
\difficulty{4}
For any counting-algorithm state $\sigma$ and natural number $z$, replacing coordinate
$9$ of its register vector by $z$ yields the register vector of the state obtained by
replacing its known-ball size $cp$ by $z$.
\end{lemma}

\begin{lemma}[Updating the outer vertex]
\lean{Complexity.LogProg.ISCount.St.upd_v}
\label{Complexity.LogProg.ISCount.St.upd_v}
\leanok
\uses{Complexity.LogProg.ISCount.St.toV, Complexity.LogProg.ISCount.rv}
\difficulty{4}
For any counting-algorithm state $\sigma$ and natural number $z$, replacing coordinate
$10$ of the register vector of $\sigma$ by $z$ gives the register vector of the state
obtained by replacing its outer-vertex register $v$ by $z$.
\end{lemma}

\begin{lemma}[Updating the next-ball count]
\lean{Complexity.LogProg.ISCount.St.upd_c}
\label{Complexity.LogProg.ISCount.St.upd_c}
\leanok
\uses{Complexity.LogProg.ISCount.St.toV, Complexity.LogProg.ISCount.rc}
\difficulty{4}
For any counting-algorithm state $\sigma$ and natural number $z$, updating its register
vector at index $11$ to $z$ gives the register vector of the state obtained by replacing
its next-ball count $c$ with $z$.
\end{lemma}

\begin{lemma}[Updating the inner vertex]
\lean{Complexity.LogProg.ISCount.St.upd_u}
\label{Complexity.LogProg.ISCount.St.upd_u}
\leanok
\uses{Complexity.LogProg.ISCount.St.toV, Complexity.LogProg.ISCount.ru}
\difficulty{4}
For any counting-algorithm state $\sigma$ and natural number $z$, replacing coordinate
$12$ of the register vector of $\sigma$ by $z$ gives the register vector of the state
obtained by replacing its inner-vertex register $u$ by $z$.
\end{lemma}

\begin{lemma}[Updating the inner count]
\lean{Complexity.LogProg.ISCount.St.upd_cnt}
\label{Complexity.LogProg.ISCount.St.upd_cnt}
\leanok
\uses{Complexity.LogProg.ISCount.St.toV, Complexity.LogProg.ISCount.rcnt}
\difficulty{4}
For any counting-algorithm state $\sigma$ and natural number $z$, replacing coordinate
$13$ of its register vector by $z$ yields the register vector of the state obtained by
replacing its inner count $cnt$ by $z$.
\end{lemma}

\begin{lemma}[Updating the guessed-code coordinate]
\lean{Complexity.LogProg.ISCount.St.upd_g}
\label{Complexity.LogProg.ISCount.St.upd_g}
\leanok
\uses{Complexity.LogProg.ISCount.St.toV, Complexity.LogProg.ISCount.rg}
\difficulty{4}
For any counting-algorithm state $\sigma$ and natural number $z$, replacing coordinate
$14$ of its register vector by $z$ yields the register vector of the state obtained by
replacing its guessed code $g$ by $z$.
\end{lemma}

\begin{lemma}[Updating the current walk code]
\lean{Complexity.LogProg.ISCount.St.upd_cur}
\label{Complexity.LogProg.ISCount.St.upd_cur}
\leanok
\uses{Complexity.LogProg.ISCount.St.toV, Complexity.LogProg.ISCount.rcur}
\difficulty{4}
For any counting-algorithm state $\sigma$ and natural number $z$, replacing coordinate
$15$ of its register vector by $z$ yields the register vector of the state obtained by
replacing its current walk code $cur$ by $z$.
\end{lemma}

\begin{lemma}[Updating the walk step count]
\lean{Complexity.LogProg.ISCount.St.upd_stp}
\label{Complexity.LogProg.ISCount.St.upd_stp}
\leanok
\uses{Complexity.LogProg.ISCount.St.toV, Complexity.LogProg.ISCount.rstp}
\difficulty{4}
For any counting-algorithm state $\sigma$ and natural number $z$, replacing coordinate
$16$ of the register vector of $\sigma$ by $z$ gives the register vector of the state
obtained by replacing its walk step count $stp$ by $z$.
\end{lemma}

\begin{definition}[Bounded register state]
\lean{Complexity.LogProg.ISCount.St.Bd}
\label{Complexity.LogProg.ISCount.St.Bd}
\leanok
\uses{Complexity.LogProg.ISCount.St}
A register state is bounded by a natural number $B$ if each of its 17 registers is at
most $B$.
\end{definition}

\begin{definition}[Phase-dependent register index]
\lean{Complexity.LogProg.ISCount.rY}
\label{Complexity.LogProg.ISCount.rY}
\leanok
\uses{Complexity.LogProg.ISCount.rlg, Complexity.LogProg.ISCount.rs}
For a Boolean phase indicator, the selected register index is the window-radius register
when the indicator is true, and the halving-count register when it is false.
\end{definition}

\begin{definition}[Boolean-dependent vertex register index]
\lean{Complexity.LogProg.ISCount.wT}
\label{Complexity.LogProg.ISCount.wT}
\leanok
\uses{Complexity.LogProg.ISCount.ru, Complexity.LogProg.ISCount.rv}
For a Boolean indicator, the selected register index is the inner-vertex register when
the indicator is true, and the outer-vertex register when it is false.
\end{definition}

\begin{definition}[Boolean-dependent stage register index]
\lean{Complexity.LogProg.ISCount.wB}
\label{Complexity.LogProg.ISCount.wB}
\leanok
\uses{Complexity.LogProg.ISCount.ri, Complexity.LogProg.ISCount.ri1}
For a Boolean indicator, the selected register index is the stage register when the
indicator is true, and the stage-plus-one register when it is false.
\end{definition}

\begin{lemma}[Boolean-selected vertex value]
\lean{Complexity.LogProg.ISCount.toV_wT}
\label{Complexity.LogProg.ISCount.toV_wT}
\leanok
\uses{Complexity.LogProg.ISCount.St, Complexity.LogProg.ISCount.St.toV, Complexity.LogProg.ISCount.wT}
\difficulty{4}
Let a counting-algorithm state have inner vertex $u$ and outer vertex $v$, and let $b$
be Boolean. The register vector entry at the index selected by $b$ is $u$ if $b$ is true
and $v$ if $b$ is false.
\end{lemma}

\begin{lemma}[Boolean-selected stage value]
\lean{Complexity.LogProg.ISCount.toV_wB}
\label{Complexity.LogProg.ISCount.toV_wB}
\leanok
\uses{Complexity.LogProg.ISCount.St, Complexity.LogProg.ISCount.St.toV, Complexity.LogProg.ISCount.wB}
\difficulty{4}
Let a counting-algorithm state have stage $i$ and stage-plus-one $i1$, and let $b$ be
Boolean. The register vector entry at the index selected by $b$ is $i$ if $b$ is true
and $i1$ if $b$ is false.
\end{lemma}

\begin{lemma}[Phase-selected register value]
\lean{Complexity.LogProg.ISCount.toV_rY}
\label{Complexity.LogProg.ISCount.toV_rY}
\leanok
\uses{Complexity.LogProg.ISCount.St, Complexity.LogProg.ISCount.St.toV, Complexity.LogProg.ISCount.rY}
\difficulty{4}
Let a counting-algorithm state have window radius $s$ and halving count $lg$, and let
$ph$ be Boolean. The register vector entry at the index selected by $ph$ is $s$ if $ph$
is true and $lg$ if $ph$ is false.
\end{lemma}

\begin{definition}[Boolean-dependent vertex value]
\lean{Complexity.LogProg.ISCount.St.wt}
\label{Complexity.LogProg.ISCount.St.wt}
\leanok
\uses{Complexity.LogProg.ISCount.St}
For a register state and a Boolean indicator, the selected vertex value is the state's
inner vertex when the indicator is true, and its outer vertex when the indicator is
false.
\end{definition}

\begin{definition}[Boolean-dependent stage value]
\lean{Complexity.LogProg.ISCount.St.wb}
\label{Complexity.LogProg.ISCount.St.wb}
\leanok
\uses{Complexity.LogProg.ISCount.St}
For a register state and a Boolean indicator, the selected stage value is the state's
stage when the indicator is true, and its stage plus one when the indicator is false.
\end{definition}

\begin{definition}[Clear machine configuration]
\lean{Complexity.LogProg.ISCount.St.Clear}
\label{Complexity.LogProg.ISCount.St.Clear}
\leanok
\uses{Complexity.LogProg.ISCount.St}
A configuration of the counting procedure is clear when its guessed code, current walk
code, walk step count, inner vertex, and inner count are all zero.
\end{definition}

\begin{definition}[Binary size with offset]
\lean{Complexity.LogProg.ISCount.lgv}
\label{Complexity.LogProg.ISCount.lgv}
\leanok
For a natural number $n$, define its size with offset to be the binary size of $n$ plus
one.
\end{definition}

\begin{definition}[Oracle reachability relation]
\lean{Complexity.LogProg.ISCount.RelO}
\label{Complexity.LogProg.ISCount.RelO}
\leanok
\uses{Complexity.LogProg.SavitchARM.Rel}
For an oracle $O:\text{Fin}(4)\to\text{List}(\text{Bool})\to\text{Bool}$, an input bit string
$x$, and a natural number $s$, define a relation on natural numbers by the reachability
relation associated with $O$, $x$, and $s$ in Savitch's alternating reachability
machine.
\end{definition}

\begin{definition}[Oracle initialization predicate]
\lean{Complexity.LogProg.ISCount.InitO}
\label{Complexity.LogProg.ISCount.InitO}
\leanok
\uses{Complexity.LogProg.SavitchARM.InitO}
Given an oracle $O:\text{Fin}(4)\to\text{List}(\text{Bool})\to\text{Bool}$, an input bit
string $x$, and a natural number $s$, define the initialization predicate on a natural
number $u$ to hold exactly when the oracle's query at index $1$, on the nested
self-delimiting encoding of $x$, the binary expansion of $u$, and the binary expansion
of $s$, returns true.
\end{definition}

\begin{definition}[Oracle acceptance predicate]
\lean{Complexity.LogProg.ISCount.AccO}
\label{Complexity.LogProg.ISCount.AccO}
\leanok
\uses{Complexity.LogProg.SavitchARM.AccO}
Given an oracle $O:\text{Fin}(4)\to\text{List}(\text{Bool})\to\text{Bool}$, an input bit
string $x$, and a natural number $s$, define the acceptance predicate on a natural
number $v$ to hold exactly when the oracle's query at index $2$, on the nested
self-delimiting encoding of $x$, the binary expansion of $v$, and the binary expansion
of $s$, returns true.
\end{definition}
